%% ma: L10-B14-0018
%% lop: 10
%% chuong: 5
%% bai: 14
%% dang: 10.14.2
%% loai: TN4
%% muc: TH
%% dap_an: "C"
%% nguon: "chuyen-de-cac-so-dac-trung-cua-mau-so-lieu-khong-ghep-nhom-toan-10-chuong-trinh-moi.pdf (GV. Phan Nhật Linh) (P04), Bài 03 Dạng 1 Phần I Câu 8"
%% kien_thuc: "Lớp 10 Bài 13 (tứ phân vị), Bài 14 (khoảng tứ phân vị)"
%% trang_thai: chinh-thuc
%% ghi_chu: "Đề gốc ghi 20 ngày nhưng chỉ liệt kê 19 số và đáp án gốc 22 (Q1 = 10, Q3 = 32) không khớp dữ liệu; đã sửa lời dẫn thành 19 ngày, đáp án đúng 24 (C). Tự duyệt theo yêu cầu giáo viên 22/10/2026."
\begin{ex}
Số ly trà sữa một quán nước bán được trong $19$ ngày qua là:
\begin{center}
\begin{tabular}{|c|c|c|c|c|c|c|c|c|c|}\hline
$4$ & $5$ & $6$ & $8$ & $9$ & $11$ & $13$ & $16$ & $18$ & $20$\\\hline
$21$ & $25$ & $30$ & $31$ & $33$ & $36$ & $37$ & $40$ & $41$ & \\\hline
\end{tabular}
\end{center}
Khoảng tứ phân vị của mẫu số liệu là
\choice{$20$}{$22$}{\True $24$}{$26$}
\loigiai{Dãy đã sắp xếp, $n=19$: $Q_2=x_{10}=20$; nửa dưới gồm $9$ giá trị đầu nên $Q_1=x_5=9$; nửa trên gồm $9$ giá trị cuối nên $Q_3=x_{15}=33$.

$\Delta_Q=33-9=24$. Chọn C.}
\end{ex}
